\chapter*{CONTINUAL LEARNING AND CATASTROPHIC FORGETTING IN LANGUAGE MODELS FOR VIETNAMESE LAW\\ABSTRACT}

Large Language Models (LLMs) reason strongly across domains, including legal question-answering and inference. Keeping their knowledge temporally valid is difficult: statutes are amended, regulations change, and confidential or repealed content must be removed. Continual Pre-Training and Continual Instruction Tuning update weights directly, so they suffer from catastrophic forgetting and weight drift, while unlearning specific facts from weights is costly and hard to verify.

This work reviews the field and proposes a RAG-centric, non-parametric framework for Vietnamese legal LLMs that keeps the base model frozen and delegates facts to an external store. A selective memory module applies Ebbinghaus-based consolidation and decay to shrink the index and reduce noise, while a compliance hard-gate guarantees that in-force clauses are never omitted. A controllable unlearning gate combining a Retriever Policy ($P$) and a Prompt Constraint ($Q$) handles temporal, access-control, and personal-data deletion, and lightweight LoRA adapters supply legal reasoning behaviour without storing facts.

The domain-general framework is validated on VLQA, ViLegalNLI, and LegalSLM, with concrete test scenarios drawn from data-rich legal sub-domains such as labour and tax law, through which knowledge update, sequential forgetting, and unlearning are measured. The approach offers compliant, low-cost, instantly updatable legal assistance with no factual weight drift.
